\documentclass[11pt,letterpaper]{article}
\usepackage[T1]{fontenc}
\usepackage[utf8]{inputenc}
\usepackage{lmodern,microtype}
\usepackage[margin=0.79in,headheight=15pt]{geometry}
\usepackage{parskip,enumitem,longtable,tabularx,booktabs,array,xcolor,fancyhdr,hyperref}
\definecolor{navy}{HTML}{17314C}
\definecolor{teal}{HTML}{087F78}
\hypersetup{colorlinks=true,linkcolor=navy,urlcolor=teal}
\setlength{\emergencystretch}{3em}
\setlist[itemize]{leftmargin=*,nosep,topsep=5pt}
\setlist[enumerate]{leftmargin=*,itemsep=3pt}
\renewcommand{\arraystretch}{1.22}
\pagestyle{fancy}\fancyhf{}
\fancyhead[L]{\small Enterprise Data Modeling}\fancyhead[R]{\small Documentation Suite}
\fancyfoot[C]{\thepage}
\newcommand{\srcbook}{\textit{The Data Model Resource Book}, revised edition, Volume 2}
\newcommand{\ResponseLine}{\par\vspace{1.05em}\noindent\rule{\linewidth}{0.4pt}\par}
\newcommand{\ResponseArea}{\ResponseLine\ResponseLine\ResponseLine}
\begin{document}
\begin{center}
{\LARGE\bfseries\color{navy} Manufacturing BOM and Revision Control}\par
\vspace{0.25em}{\large Practical modeling laboratory}\par
\vspace{0.3em}{\small Original engineering study material \textbullet{} 2026 edition}
\end{center}
\vspace{0.8em}
\section{Scenario and goal}
An assembly plant with revision-controlled parts, substitute components and production runs

This is an original lab. Use the chapter references and industry model sheets for terminology, but all example facts are hypothetical. Time-box discovery to 30 minutes, drafting to 60 minutes, and review and tests to 45 minutes.


\section{Inputs and tools}
\begin{itemize}
\item Input: a one-page business narrative, actor list and at least six example records.
\item Tools: pencil or diagram-as-code editor, PostgreSQL where SQL is used, and model review checklists.
\item Required artifacts: glossary, conceptual ERD, physical relation inventory, test plan and rationale.
\end{itemize}

\section{Procedure}
\begin{enumerate}
\item Design Part, PartRevision, BOMLine, WorkOrder and Consumption.
\item Represent part-effective dates and a quantity per assembly.
\item Explain how BOM cycles are detected transactionally.
\item Separate engineering BOM from as-built inventory.
\end{enumerate}

\section{Adversarial test data}
\begin{itemize}
\item A part substitutes for itself.
\item Two active revision ranges overlap.
\item A completed run consumes more than stocked quantity.
\end{itemize}
Classify each test as schema-enforced, transactionally enforced or policy-enforced, and explain the enforcement boundary.


\section{Verification and grading}
\begin{tabularx}{\linewidth}{@{}p{0.25\linewidth} p{0.71\linewidth}@{}}
\toprule
\textbf{Dimension} & \textbf{Check} \\ 
\midrule
Semantics & Each table and relation has an unambiguous domain definition \\ 
Keys & All identities, foreign keys and uniqueness scopes documented \\ 
History & State transitions, effective time and correction process unambiguous \\ 
Testability & Positive and negative cases produce a specified observable outcome \\ 
Review & An independent reviewer can trace requirements to constraints \\ 
\bottomrule
\end{tabularx}

\section{Submission}
Deliver the diagram source, rendered view, table specification, test evidence, unresolved-question register and an annotated scorecard. Do not state that a model is production-ready solely because a diagram renders.


\vfill
\hrule
{\footnotesize\textbf{Provenance.} Supplementary engineering practice: original content; the cited book provides domain-modeling context.}
\end{document}
